\documentclass[10pt]{extarticle}

% ---------- Packages ----------
\usepackage{fancyhdr}
\usepackage{geometry}
\usepackage{hyperref}
\usepackage{xcolor}
\usepackage{amsmath, amsfonts, bm}
\usepackage{subcaption}
\usepackage{graphicx}
\usepackage{tikz}
\usepackage{nomencl}
\usepackage{array}
\usepackage{longtable}
\usepackage{pdflscape}
\usepackage[export]{adjustbox}
\usepackage[utf8]{inputenc}
\usepackage{parskip}

\usetikzlibrary{shapes.geometric}
\usetikzlibrary{arrows.meta, arrows}

% ---------- Page setup ----------
\geometry{
    a4paper,
    total={170mm,257mm},
    left=20mm,
    top=20mm
}

% ---------- Headers and footers ----------
\pagestyle{fancy}
\pagenumbering{arabic}

\fancyhead[L]{Rocket Physics by Yaroslav Knyazkov -- v2 draft}
\fancyfoot[C]{}
\fancyfoot[R]{\thepage}

% ---------- Hyperlinks ----------
\hypersetup{
    colorlinks=true,
    linkcolor={blue!50!black},
    citecolor={blue!50!black},
    urlcolor={blue!80!black}
}

% ---------- Equation numbering ----------
\numberwithin{equation}{subsection}

% ---------- Nomenclature ----------
% \nomenclature{symbol}{description} calls live next to each symbol's definition in the section
% files; \printnomenclature below collects and prints them. Needs one extra build step beyond
% pdflatex/latexmk's default pass: run `makeindex main.nlo -s nomencl.ist -o main.nls` once (or
% configure your latexmk to do it automatically), then recompile -- otherwise the list prints
% empty/stale. No sort keys or \nomgroup categories (Roman/Greek/Abbreviations) are set up yet --
% entries sort by makeindex's default, which can look odd for \vec{}-wrapped symbols; add
% \nomenclature[sortkey]{...}{...} keys later if that becomes annoying.
\makenomenclature

% ---------- Custom commands (carried over from v1; drop any you don't need) ----------
\newcommand{\todo}[1]{%
    \textcolor{red}{(#1)}%
}
% Two kinds of to-do, visually distinct from each other and from the blue AI text:
%   \simtodo{...}  implemented in the paper, not yet in the simulator (dark red, "sim:")
%   \restodo{...}  needs research first (orange, "research:")
\definecolor{simtodocolor}{RGB}{170,0,0}
\definecolor{restodocolor}{RGB}{230,120,0}
\newcommand{\simtodo}[1]{\textcolor{simtodocolor}{(sim: #1)}}
\newcommand{\restodo}[1]{\textcolor{restodocolor}{(research: #1)}}
% \figtodo{...} marks a planned figure: a framed teal box stating the type, content and source.
\newcommand{\figtodo}[1]{%
    \par\medskip\noindent
    \fcolorbox{teal}{teal!8}{\parbox{\dimexpr\linewidth-2\fboxsep-2\fboxrule\relax}{%
        \raggedright\small\color{teal!60!black}\textbf{Figure (planned):} #1}}%
    \par\medskip}
% \comptodo{...} marks a planned model comparison: a framed grey-violet box stating the quantities
% compared on the demo rocket; the results are generated by the simulator.
\definecolor{compcolor}{RGB}{90,60,140}
\newcommand{\comptodo}[1]{%
    \par\medskip\noindent
    \fcolorbox{compcolor}{compcolor!8}{\parbox{\dimexpr\linewidth-2\fboxsep-2\fboxrule\relax}{%
        \raggedright\small\color{compcolor}\textbf{Comparison (planned):} #1}}%
    \par\medskip}
% AI-written text is blue; the author's own text stays black. The author reviews the document
% part by part and turns a reviewed part black by removing its wrapper.
%   \ai{...}               inline or multi-paragraph (\newcommand is \long by default)
%   \begin{aitext}...\end{aitext}   whole sections / \input files
% \todo stays red inside either: its own \textcolor overrides the surrounding blue.
\newcommand{\ai}[1]{{\color{blue}#1}}
\newenvironment{aitext}{\color{blue}}{}

\begin{document}

\tableofcontents

\newpage
\printnomenclature

% ---------- Section scaffold (v2 draft) ----------
% Mirrors v1's section breakdown as a starting point only -- rename, merge, split or drop freely.
% Each \input target below is an empty stub in sections/; fill them in, or delete the \input line
% and the stub file if a section does not apply to v2.

\newpage
\section{Systems}
\subsection{Flight}
Flight of the rocket can be divided into an atmospheric and a space part. The equations of motion
for the two are the same; only the magnitudes of the forces differ.

\subsubsection{Forces \& Moments - Atmospheric flight}
\paragraph{Weight} is the main force to overcome when lifting something off the ground. It acts at
the centre of gravity (CoG), not the centre of mass (CoM); the two coincide unless gravity varies
appreciably over the body, an effect called the gravity gradient, which only matters in space. For
atmospheric flight the two points can therefore usually be assumed to be one and the same. The
weight force acts in the direction of the gravity field $\hat{z}_g$ and has magnitude $mg$. It can
therefore be written as
\begin{equation}
    \vec{W} = mg\hat{z}_g.
\end{equation}
\nomenclature{CoG}{centre of gravity}
\nomenclature{CoM}{centre of mass}
\nomenclature{$m$}{total mass of the rocket}
\nomenclature{$g$}{gravitational acceleration}
\nomenclature{$\hat{z}_g$}{direction of the local gravity field}
\nomenclature{$\vec{W}$}{weight force}
More on these variables in \autoref{sec:mod_mass} (mass) and \autoref{sec:mod_gravity} (gravity).
% total mass

% gravity models - direction and magnitude

\paragraph{Thrust} is what overcomes the weight, by ``brute force''. Thrust comes from ejecting
propellant at high velocity, plus the pressure difference across the nozzle exit. It acts at the
centre of thrust (CoT) and is usually given as a thrust magnitude (over time) $T$ and a thrust
direction $\hat{x}_t$, either constant or varying over time. The thrust force is given as
\begin{equation}
    \vec{T} = T\hat{x}_t.
\end{equation}
\nomenclature{CoT}{centre of thrust}
\nomenclature{$T$}{thrust magnitude}
\nomenclature{$\hat{x}_t$}{thrust direction unit vector}
\nomenclature{$\vec{T}$}{thrust force vector}
It is worth noting that a rocket can have multiple thrust application points, e.g.\ for higher
thrust or for vehicle attitude control. These give an equivalent total force and moment:
\begin{equation}
    \vec{T}_{t, f} = \sum_i \vec{T}_i,\quad \vec{T}_{t, m}(a) = \sum_i \vec{aT_i}\times\vec{T}_i.
\end{equation}
\nomenclature{$\vec{T}_i$}{thrust force vector at the $i$-th application point}
\nomenclature{$\vec{T}_{t,f}$}{total thrust force, summed over application points}
\nomenclature{$\vec{T}_{t,m}(a)$}{total thrust moment about reference point $a$, summed over application points}
More on these variables in \autoref{sec:mod_thrust}.
% linintrp
% pressure variable
% liquid motor
% solid motor

\paragraph{Aerodynamics} is the outcome of flying within the atmosphere. Aerodynamic loads arise
from the pressure and shear-stress distributions over the airframe's surface. The force components
along the rocket's own axes are called drag, lift and side force; by convention each acts opposite
to the (positive) axis it is projected onto. The net force acts at the centre of pressure (CoP),
about which the moment of these forces is zero by definition. In practice it is more convenient to
use a fixed reference point instead of the true, state-dependent CoP, and compute moments about
that reference point instead. The force and moment can be written as
\begin{equation}
    \vec{A}_f = -A_{x,f}\hat{x}_r -A_{y,f}\hat{y}_r -A_{z,f}\hat{z}_r,\quad
    \vec{A}_m = -A_{x,m}\hat{x}_r -A_{y,m}\hat{y}_r -A_{z,m}\hat{z}_r.
\end{equation}
\nomenclature{CoP}{centre of pressure}
\nomenclature{$\hat{x}_r,\hat{y}_r,\hat{z}_r$}{rocket body-frame unit vectors (forward, right, down)}
\nomenclature{$\vec{A}_f$}{aerodynamic force vector}
\nomenclature{$\vec{A}_m$}{aerodynamic moment vector}
\nomenclature{$A_{x,f},A_{y,f},A_{z,f}$}{aerodynamic force components along the body axes}
\nomenclature{$A_{x,m},A_{y,m},A_{z,m}$}{aerodynamic moment components along the body axes}
More on these force magnitudes in \autoref{sec:mod_aerodynamics}.
% Derivation of reference center of pressure
% Variables - Constant, Linearization,
% Wing body

\subsubsection{Forces - Space flight}
\todo{Add }
\subsubsection{Equations of motion - Point mass}
\subsubsection{Equations of motion - 2D}
\subsubsection{Equations of motion - 3D}


% ---------- Bibliography ----------
% Left out for now -- v1's full bibliography lives in ../latex/main.tex (lines ~127-186).
% Copy over only the \bibitem entries v2 actually cites, when it cites them.
% \newpage
% \begin{thebibliography}{99}
% \small
% \end{thebibliography}

\end{document}
